% !PRINTED 579
\weberpage{579}

%<<CONT
Dies widerspricht aber dem Satze 5. des vorigen Paragraphen, nach dem $\tau$ nach einem Primzahlmodul $p$ keiner Congruenz von niedrigerem als $n^{\mathrm{tem}}$ Grade genügen kann.

Demnach ergiebt sich aus (3), \gesperrt{dass die Grundzahl $\Delta$ des Körpers, vom Vorzeichen abgesehen, die absolute Norm der Function $F'(\tau)$ ist}, dass also
\weq{5}{\pm \Delta = N_a[F'(\tau)]}
zu setzen ist.

Wenn nun statt der $\omega_i$ eine andere Basis $\omega_i'$ von $\mathfrak{o}$ zu Grunde gelegt wird, so tritt an Stelle von $\tau$ eine andere Form
\weq{6}{\tau' = \omega_1' t_1 + \omega_2' t_2 + \cdots + \omega_n' t_n,}
wofür wir auch setzen können
\weq{7}{\tau' = \omega_1 t_1' + \omega_2 t_2' + \cdots + \omega_n t_n',}
und darin sind die $\omega_i'$ mit den $\omega_i$ durch eine lineare Substitution mit der Determinante $\pm 1$ verbunden (\S.\,145):
\weq{8}{(\omega_1', \omega_2', \ldots, \omega_n') = C(\omega_1, \omega_2, \ldots, \omega_n).}
Führt man diese Substitution in (6) aus, und ordnet nach $\omega_1, \omega_2, \ldots, \omega_n$, so ergiebt sich
\weq{9}{(t_1', t_2', \ldots, t_n') = C_1(t_1, t_2, \ldots, t_n),}
wenn $C_1$ die transponirte Substitution von $C$ ist (\S.\,37). Bildet man nun die Function $F(t)$ für die Function $\tau'$, so mag sich $F_1(t)$ ergeben; die Ableitung für $t=\tau'$ sei $F_1'(\tau')$. Setzen wir nun
\[
F'(\tau) = \Psi(t_1, t_2, \ldots, t_n),
\]
so ist $\Psi$ eine ganze Function in $R$, und es ergiebt sich wegen (7)
\[
F_1'(\tau') = \Psi(t_1', t_2', \ldots, t_n').
\]

% [sic: The source prints $F'(\tau')$ as the first functional; the preceding definitions compare $F'(\tau)=\Psi(t_1,\ldots,t_n)$ with $F_1'(\tau')=\Psi(t_1',\ldots,t_n')$. The primed argument is retained as printed.]
Da die Substitution (9) umkehrbar ist, so folgt hieraus, dass die beiden Functionale $F'(\tau')$ und $F_1'(\tau')$ gegenseitig durch einander theilbar sind (\S.\,143), und dass sie mithin associirt sind.

Die Function $F'(\tau)$, deren absolute Norm gleich dem absoluten Werthe der Grundzahl ist, nennen wir daher das \gesperrt{Grundfunctional}, und das durch $F'(\tau)$ repräsentirte Ideal das \gesperrt{Grundideal} des Körpers $\Omega$. Nun sei $\mathfrak{p}^e$ die höchste Potenz des Primideals $\mathfrak{p}$, die in $p$ aufgeht, und $e \geqq 1$, dann können wir nach 5. des vorigen Paragraphen
\weq{10}{F(t) \equiv P(t)^e\, \Phi(t)\ (\mathrm{mod}\ p)}
setzen, worin $\Phi(t)$ eine ganze Function in $R$ von der Beschaffenheit ist, dass $\Phi(\tau)$ nicht durch $\mathfrak{p}$ theilbar ist. Aus (10) aber
%CONT>>
